\subsection{Aggregate decomposition across countries}

Let \(q_{c,t}\) denote country \(c\)'s share of employees in the covered-country
sample and let \(p_{ce,t}\) denote education group \(e\)'s share of employees
within that country. The sample-wide mean wage is
\begin{equation}
  R_t^{G}
  =
  \sum_c q_{c,t} R_{c,t},
  \qquad
  R_{c,t}
  =
  \sum_e p_{ce,t} r_{ce,t}.
  \label{eq:global_wage_identity}
\end{equation}
Applying the symmetric decomposition first across countries and then within
each country gives
\begin{align}
  R_1^{G}-R_0^{G}
  ={}&
  \sum_c \bar{R}_c\left(q_{c,1}-q_{c,0}\right)
  \nonumber\\
  &+
  \sum_c \bar{q}_c
  \sum_e \bar{r}_{ce}\left(p_{ce,1}-p_{ce,0}\right)
  \nonumber\\
  &+
  \sum_c \bar{q}_c
  \sum_e \bar{p}_{ce}\left(r_{ce,1}-r_{ce,0}\right),
  \label{eq:global_symmetric_decomposition}
\end{align}
where a bar denotes the arithmetic mean of the two endpoints. The first term
is the between-country composition component. The second is educational
upgrading within countries, and the third is the wage change within
country--education cells. This hierarchical decomposition is exactly
additive. The between-country component should not be interpreted as worker
migration because country weights also change with population, labor-force
participation, employee status, and survey coverage.

The baseline aggregate fixes country weights at
\begin{equation}
  q_c^{*}
  =
  \frac{N_{c,0}+N_{c,1}}
       {\sum_j\left(N_{j,0}+N_{j,1}\right)}.
  \label{eq:fixed_country_weights}
\end{equation}
With \(q_c^{*}\) held fixed, the between-country component is zero and the
sample-wide change separates into the original two components: educational
upgrading and wage changes within education groups. Equal weights by country
and changing employment weights are reported as robustness exercises.
